\documentclass[11pt]{article}
\usepackage{fontspec}
\usepackage{polyglossia}
\setmainlanguage{russian}
\setotherlanguage{english}
\setmainfont{FreeSerif}
\setmonofont{FreeMono}
\usepackage{graphicx}
\usepackage{amsmath}
\usepackage{amssymb}
\usepackage{booktabs}
\usepackage{geometry}
\usepackage{hyperref}
\usepackage{microtype}
\usepackage{caption}
\geometry{margin=2.5cm}

\title{Сравнение стратегий retrieval для RAG-системы \\ по документации 1С:Предприятие}
\author{Денисенко Алексей Сергеевич }
\date{Июль 2026}

\begin{document}
\maketitle

\begin{abstract}
Работа посвящена пилотному исследованию компонента retrieval в системах
Retrieval-Augmented Generation (RAG), применяемых для ответов на вопросы
разработчиков по платформе 1С:Предприятие. Мы сравниваем три стратегии
поиска релевантных фрагментов знаний: классический лексический baseline
на основе TF-IDF, латентно-семантические (LSA) эмбеддинги, обученные
локально на корпусе методом усечённого SVD, и гибридный метод на основе
Reciprocal Rank Fusion (RRF). Эксперимент проведён на компактном
авторском корпусе из 20 карточек знаний по объектам метаданных 1С и
25 вопросах с разметкой релевантности. Результаты показывают, что LSA
превосходит TF-IDF по MRR и Recall@1, при этом гибридный метод обеспечивает
наиболее устойчивое ранжирование на верхних позициях. Обсуждаются
ограничения подхода и направления дальнейшей работы, включая масштабирование
корпуса и добавление компонента генерации ответа.
\end{abstract}

\section{Введение}

Платформа 1С:Предприятие широко используется в корпоративной автоматизации
на русскоязычном рынке, а разработка на встроенном языке платформы требует
постоянного обращения к обширной документации, стандартам разработки и базе
знаний сообщества. Ручной поиск релевантной информации в этой документации
занимает значительное время разработчика. Системы Retrieval-Augmented
Generation (RAG) позволяют автоматизировать этот процесс: на основе вопроса
пользователя система находит релевантные фрагменты базы знаний и использует
их как контекст для генерации ответа языковой моделью.

Качество итогового ответа RAG-системы напрямую определяется качеством
компонента retrieval: если релевантный фрагмент не найден на этапе поиска,
языковая модель не получит нужный контекст и не сможет дать точный ответ.
Поэтому в данной работе мы сосредоточены именно на сравнении стратегий
retrieval, а не на компоненте генерации.

\textbf{Вклад работы.} (1) Собран компактный авторский корпус карточек
знаний по ключевым объектам платформы 1С:Предприятие и QA-датасет с
разметкой релевантности для оценки retrieval. (2) Реализованы и сравнены
три стратегии поиска: TF-IDF, LSA-эмбеддинги, обученные локально без
использования внешних предобученных моделей, и их гибрид на основе RRF.
(3) Проведён количественный анализ по метрикам Recall@k и MRR, а также
анализ задержки (latency) каждого метода.

\section{Обзор литературы}

Классические методы информационного поиска, такие как TF-IDF и BM25,
остаются сильным baseline для задач поиска релевантных документов
благодаря простоте, интерпретируемости и отсутствию необходимости в
обучающих данных. Латентно-семантический анализ (LSA), основанный на
сингулярном разложении (SVD) матрицы term-document, позволяет получить
низкоразмерные представления документов, учитывающие скрытые семантические
связи между терминами, и является одним из первых подходов к построению
плотных представлений текста, предшествующим современным нейросетевым
эмбеддингам.

Современные RAG-системы, как правило, используют предобученные нейросетевые
эмбеддинги (например, модели семейства sentence-transformers) для плотного
retrieval, что даёт заметный прирост по сравнению с чисто лексическими
методами на больших и разнородных корпусах. Гибридные подходы, комбинирующие
лексический и плотный поиск через методы объединения ранжирований, такие
как Reciprocal Rank Fusion (RRF), показывают устойчивое улучшение полноты
поиска по сравнению с использованием одного метода, поскольку лексический
и семантический поиск компенсируют слабые стороны друг друга: лексический
поиск лучше находит точные термины и аббревиатуры, а семантический --
перефразированные запросы.

Применительно к предметной области 1С:Предприятие специализированные
исследования retrieval для RAG-систем на момент подготовки работы
практически отсутствуют в открытом доступе, что делает данную предметную
область подходящим полигоном для пилотного исследования. Отличие настоящей
работы от типичных промышленных RAG-пайплайнов заключается в отказе от
внешних предобученных эмбеддинг-моделей: LSA-представления обучаются
исключительно на локальном корпусе, что позволяет оценить, насколько
эффективны легковесные локально обучаемые методы в условиях ограниченных
вычислительных ресурсов и отсутствия доступа к внешним API.

\section{Метод}

\subsection{Постановка задачи}

Пусть $D = \{d_1, \dots, d_n\}$ -- корпус документов (карточек знаний), а
$q$ -- вопрос пользователя. Задача retrieval заключается в построении
ранжирования $R(q) = (d_{i_1}, d_{i_2}, \dots, d_{i_n})$ документов по
убыванию релевантности вопросу $q$, такого что документы, содержащие
ответ на $q$, оказываются на верхних позициях ранжирования.

\subsection{TF-IDF baseline}

Каждый документ $d_j$ и запрос $q$ представляются в виде векторов TF-IDF
с учётом униграмм и биграмм:
\begin{equation}
\text{tfidf}(t, d) = \text{tf}(t, d) \cdot \log\frac{N}{\text{df}(t)},
\end{equation}
где $\text{tf}(t, d)$ -- частота термина $t$ в документе $d$, $\text{df}(t)$ --
число документов, содержащих термин $t$, $N$ -- общее число документов.
Ранжирование документов производится по убыванию косинусной близости между
вектором запроса и векторами документов.

\subsection{LSA-эмбеддинги}

Матрица TF-IDF $X \in \mathbb{R}^{n \times m}$ подвергается усечённому
сингулярному разложению:
\begin{equation}
X \approx U_k \Sigma_k V_k^\top,
\end{equation}
где $k=10$ -- размерность латентного пространства, подобранная вручную с
учётом малого объёма корпуса. Строки $U_k \Sigma_k$ используются как плотные
представления документов; запрос проецируется в то же пространство через
обученную матрицу $V_k$. Ранжирование строится по косинусной близости в
латентном пространстве.

\subsection{Гибридный метод (RRF)}

Ранжирования, полученные TF-IDF и LSA, объединяются методом Reciprocal
Rank Fusion:
\begin{equation}
\text{RRF}(d) = \sum_{r \in \{R_{\text{tfidf}}, R_{\text{lsa}}\}} \frac{1}{k + \text{rank}_r(d)},
\end{equation}
где $\text{rank}_r(d)$ -- позиция документа $d$ в ранжировании $r$ (начиная
с 1), $k=60$ -- константа сглаживания, стандартное значение, используемое в
литературе по RRF.

\section{Эксперименты}

\subsection{Данные}

Для эксперимента собран авторский корпус из $n=20$ карточек знаний,
описывающих ключевые объекты метаданных и языковые механизмы платформы
1С:Предприятие (справочники, документы, регистры сведений и накопления,
язык запросов, встроенный язык, обработчики событий, роли и права доступа
и др.). На основе корпуса составлен QA-датасет из 25 вопросов, каждый из
которых сформулирован с использованием синонимов и перефразирований
относительно исходного текста документов, чтобы избежать тривиального
лексического совпадения и проверить устойчивость методов к парафразу.
Каждому вопросу сопоставлен один или два релевантных документа (gold labels).

\subsection{Экспериментальная настройка}

Все три метода реализованы на языке Python с использованием библиотеки
scikit-learn (\texttt{TfidfVectorizer}, \texttt{TruncatedSVD}). Стоп-слова
русского языка удалены на этапе токенизации. Случайное состояние (random
seed) зафиксировано равным 42 для полной воспроизводимости результатов.
Эксперимент выполнялся на CPU без использования GPU; ввиду небольшого
объёма корпуса и малого числа компонент SVD вычисления занимают менее
секунды.

\subsection{Метрики}

Качество ранжирования оценивается метриками Recall@1, Recall@3, Recall@5
(доля вопросов, для которых хотя бы один релевантный документ попал в
top-$k$ ранжирования) и Mean Reciprocal Rank (MRR), учитывающей позицию
первого релевантного документа в ранжировании. Дополнительно измерена
средняя задержка (latency) обработки одного запроса.

\subsection{Результаты}

Таблица~\ref{tab:results} и Рисунок~\ref{fig:recall} представляют результаты
сравнения трёх методов.

\begin{table}[h]
\centering
\caption{Сравнение методов retrieval по метрикам Recall@k, MRR и задержке}
\label{tab:results}
\begin{tabular}{lccccc}
\toprule
Метод & Recall@1 & Recall@3 & Recall@5 & MRR & Latency (мс) \\
\midrule
TF-IDF (baseline) & 0.60 & 0.80 & 0.88 & 0.723 & 0.63 \\
LSA (SVD-эмбеддинги) & \textbf{0.68} & 0.76 & 0.84 & \textbf{0.744} & \textbf{0.38} \\
Гибрид (RRF) & 0.64 & \textbf{0.80} & 0.84 & 0.731 & 1.02 \\
\bottomrule
\end{tabular}
\end{table}

\begin{figure}[h]
\centering
\includegraphics[width=0.75\textwidth]{../results/recall_comparison.png}
\caption{Сравнение Recall@k для трёх методов retrieval}
\label{fig:recall}
\end{figure}

LSA-эмбеддинги показали наилучший результат по MRR (0.744) и Recall@1
(0.68), превзойдя TF-IDF baseline (MRR 0.723, Recall@1 0.60), что говорит о
том, что латентно-семантическое представление лучше справляется с
перефразированными вопросами за счёт учёта скрытых связей между терминами,
а не только точного совпадения слов. При этом на позициях top-3 и top-5
TF-IDF и гибридный метод показали более высокий Recall, чем чистый LSA,
что указывает на то, что чисто лексический поиск надёжнее находит документ
среди более широкого списка кандидатов, если он содержит специфические
термины из запроса. Гибридный метод на основе RRF показал наиболее
сбалансированный результат по Recall@3 (0.80, лучший результат среди всех
методов) при небольшом увеличении задержки, что подтверждает гипотезу о
взаимной компенсации сильных сторон лексического и семантического поиска.
Задержка всех методов пренебрежимо мала (менее 1.1 мс на запрос) на данном
объёме корпуса.

\section{Ограничения (Limitations)}

Результаты работы имеют ряд существенных ограничений, которые необходимо
учитывать при интерпретации.

\textbf{Размер корпуса и датасета.} Корпус из 20 документов и 25 вопросов
является пилотным и синтетическим; такой объём недостаточен для получения
статистически значимых выводов и не отражает масштаб реальной документации
1С (десятки тысяч страниц технической документации и форумных обсуждений).
Полученные значения метрик следует интерпретировать как предварительную
оценку тенденций, а не как окончательный результат.

\textbf{Отсутствие компонента генерации.} Работа сосредоточена
исключительно на этапе retrieval; полноценный RAG-пайплайн должен включать
этап генерации ответа языковой моделью на основе найденного контекста, что
не реализовано в данной версии работы из-за ограничений доступа к внешним
предобученным моделям и API в вычислительной среде на момент проведения
экспериментов. Оценка итогового качества ответа (а не только релевантности
найденного фрагмента) остаётся направлением дальнейшей работы.

\textbf{Отказ от нейросетевых эмбеддингов.} В качестве плотного
представления использован LSA, а не современные предобученные
sentence-эмбеддинги (например, модели семейства sentence-transformers или
E5), что могло привести к недооценке потенциала семантического поиска;
сравнение с такими моделями является естественным следующим шагом.

\textbf{Единственный прогон эксперимента.} Ввиду детерминированности всех
используемых методов относительно фиксированного seed эксперимент выполнен
в единственном прогоне без усреднения по нескольким seed; при увеличении
корпуса и введении стохастических компонентов (например, обучаемых
нейросетевых эмбеддингов) необходимо проводить 5+ повторов с разными seed и
приводить среднее значение со стандартным отклонением.

\textbf{Отсутствие сравнения с BM25 и продакшн-системами.} В качестве
лексического baseline использован TF-IDF с косинусной близостью, а не
классический BM25, который обычно показывает несколько лучшие результаты
на реальных корпусах; также не проведено сравнение с промышленными
RAG-фреймворками.

\section{Заключение}

В работе проведено пилотное сравнение трёх стратегий retrieval -- TF-IDF,
LSA-эмбеддингов и их гибрида на основе RRF -- для задачи поиска релевантных
фрагментов документации 1С:Предприятие в контексте RAG-систем. Результаты
показывают, что локально обученные LSA-эмбеддинги способны превзойти
лексический baseline по точности первого результата (Recall@1, MRR), а
гибридный метод обеспечивает наиболее сбалансированное качество на
top-3 позициях. Работа демонстрирует, что даже без доступа к внешним
предобученным моделям возможно построить работоспособный компонент
retrieval для доменно-специфичных RAG-систем. Направления дальнейшей
работы включают масштабирование корпуса до объёма реальной документации,
добавление компонента генерации ответа, сравнение с BM25 и современными
нейросетевыми эмбеддингами, а также проведение многократных прогонов
эксперимента для получения статистически надёжных оценок.

\begin{thebibliography}{9}

\bibitem{lewis2020}
Lewis, P., et al. Retrieval-Augmented Generation for Knowledge-Intensive NLP
Tasks. \textit{NeurIPS}, 2020.

\bibitem{deerwester1990}
Deerwester, S., Dumais, S. T., Furnas, G. W., Landauer, T. K., Harshman, R.
Indexing by Latent Semantic Analysis. \textit{Journal of the American
Society for Information Science}, 41(6), 1990.

\bibitem{robertson2009}
Robertson, S., Zaragoza, H. The Probabilistic Relevance Framework: BM25 and
Beyond. \textit{Foundations and Trends in Information Retrieval}, 3(4), 2009.

\bibitem{cormack2009}
Cormack, G. V., Clarke, C. L. A., Buettcher, S. Reciprocal Rank Fusion
Outperforms Condorcet and Individual Rank Learning Methods. \textit{SIGIR},
2009.

\end{thebibliography}

\end{document}
